\documentclass[11pt,a4paper]{article}
\usepackage{gtu-el-yazisi}
\begin{document}
\notbilgi{Sayısal İşaret İşleme}{DTFT, DFT, dairesel konvülasyon ve FFT}{Dsp-ders-13.pdf}

\sayfa{1}
\section{DTFT}
\[ x[n] \longrightarrow X(e^{j\omega}) = \sum_{n=-\infty}^{\infty} x[n]\, e^{-j\omega n} \]
\elcizimi[0.45]{sekiller/Dsp-ders-13-s1-1.png}[1]

\section{DFT}
\[ x[n] \longleftrightarrow X[k] = \sum_{n=0}^{N-1} x[n]\, e^{-j\left(\frac{2\pi}{N}k\right)n} \]
\vurgu[turuncu]{{$j\omega$} örneklendi}

\elcizimi[0.7]{sekiller/Dsp-ders-13-s1-2.png}[1]
\elcizimi[0.75]{sekiller/Dsp-ders-13-s1-3.png}[1]
\elcizimi[0.6]{sekiller/Dsp-ders-13-s1-4.png}[1]
\elcizimi[0.55]{sekiller/Dsp-ders-13-s1-5.png}[1]

\vurgu[pembe]{yeni bir konu sonra ilişkilendireceğiz}

\section{Periodic sequences}
\okunamadi{periyodik, 1 dönem}\ $\tilde{x}[n] = \tilde{x}[n+N]$ \quad tüm $n$ ler için sağlanıyorsa periyodiktir

\section{Discrete fourier series}
\[ \tilde{X}[k] = \sum_{n=0}^{N-1} \tilde{x}[n]\, e^{-j\frac{2\pi}{N}kn} \]
\[ \tilde{x}[n] = \frac{1}{N}\sum_{n=0}^{N-1} \tilde{X}[n]\, e^{j\frac{2\pi}{N}kn}\footnotemark \]
\footnotetext{Claude notu: muhtemelen toplam indisi $k$ ve $\tilde{X}[k]$ olmalı.}
\elcizimi[0.75]{sekiller/Dsp-ders-13-s1-6.png}[1]

\section{Periodic Convolution}
$\tilde{x}_1[n]$, $\tilde{x}_2[n]$ with period $N$

DFS $\rightarrow \tilde{X}_1[k]$, $\tilde{X}_2[k]$
\[ \tilde{X}_3[k] = \tilde{X}_1[k]\,\tilde{X}_2[k] \]
\[ \tilde{x}_3[n] = \sum_{m=0}^{N-1} \tilde{x}_1[m]\,\tilde{x}_2[n-m] \]
periodic convolution with period $N$ \quad (Dairesel konvülasyon)

\sayfa{2}
\section{Linear Convolution}
\[ \underset{\substack{\uparrow\\ x_3}}{y[n]} = \sum_{k=-\infty}^{\infty} \underset{\substack{\uparrow\\ x_1}}{x[k]}\; \underset{\substack{\uparrow\\ x_2}}{h[n-k]} \]
\formulkutu{x_3[n] = \sum_{m=-\infty}^{\infty} x_1[m]\, x_2[n-m]}
\[ \tilde{X}_3[k] = \mathrm{DFS}\bigl[\tilde{x}_1[n]\bigr] \cdot \mathrm{DFS}\bigl[\tilde{x}_2[n]\bigr] \]
\[ \tilde{X}_3[k] = \tilde{X}_1[k] \cdot \tilde{X}_2[k] \]
\[ x_3[n] = \mathrm{IFT}\Bigl(\underset{\substack{\uparrow\\ \text{bilgisayarda}}}{\mathrm{DFT}(x_1[n])}\ \underset{\substack{\uparrow\\ \text{bilgisayarda}}}{\mathrm{DFT}(x_2[n])}\Bigr) \]
\elcizimi[0.85]{sekiller/Dsp-ders-13-s2-1.png}[2]
\vurgu[pembe]{Bunun DFS'i ile bunun DFT'si arasında herhangi bir fark yok}

DFT'nin formülü mecburen geldi o mecburiyet bizi DFS'e yöneltti. çıkan sonuç yine dairesel konvülasyon sonucudur.

\[ \sin(2\pi\, 5t) = \sin(10\pi t) \]
$F_s =$ \qty{100}{\hertz}
\[ x[n] = \sin\!\left(\frac{\pi}{10}\,n\right) \]
\elcizimi[0.75]{sekiller/Dsp-ders-13-s2-2.png}[2]
\[ \omega_k = \frac{2\pi}{N}\,k \qquad\qquad x[n] = \sin\!\left(\frac{\pi}{10}\,n\right) \]

\sayfa{3}
\[ x[n] = \cos(0{,}1\pi n) + \cos(0{,}2\pi n) \]
\elcizimi[0.55]{sekiller/Dsp-ders-13-s3-1.png}[3]
\elcizimi[0.55]{sekiller/Dsp-ders-13-s3-2.png}[3]
\[ N=10 \;\Rightarrow\; \Delta\omega = \frac{2\pi}{10} = 0{,}2\pi \]
\[ N=20 \;\longrightarrow\; \Delta\omega = \frac{2\pi}{20} = 0{,}1\pi \]
\elcizimi[0.5]{sekiller/Dsp-ders-13-s3-3.png}[3]
\elcizimi[0.5]{sekiller/Dsp-ders-13-s3-4.png}[3]
\elcizimi[0.85]{sekiller/Dsp-ders-13-s3-5.png}[3]

Bu işaretin Fourieri alınmış olur
\elcizimi[0.6]{sekiller/Dsp-ders-13-s3-6.png}[3]

\sayfa{4}
\elcizimi[0.6]{sekiller/Dsp-ders-13-s4-1.png}[4]
\begin{align*}
X[k] &= \sum_{n=0}^{N-1} x[n]\, e^{-j\frac{2\pi}{N}kn} \\
X[k] &= \sum_{n=0}^{3} x[n]\, e^{-j\frac{2\pi}{4}kn} && k=0,1,2,3 \\
X[k] &= \sum_{n=0}^{4} x[n]\, e^{-j\frac{2\pi}{5}kn} && k=0,1,2,3,4 \\
X[k] &= \sum_{n=0}^{5} x[n]\, e^{-j\frac{2\pi}{6}kn} && k=0,1,2,3,4,5
\end{align*}
\elcizimi[0.5]{sekiller/Dsp-ders-13-s4-2.png}[4]
\elcizimi[0.55]{sekiller/Dsp-ders-13-s4-3.png}[4]
\elcizimi[0.55]{sekiller/Dsp-ders-13-s4-4.png}[4]
\elcizimi[0.55]{sekiller/Dsp-ders-13-s4-5.png}[4]

\sayfa{5}
\section{Hesap yükünü azaltalım}
\begin{align*}
X[k] &= \sum_{n=0}^{N-1} x[n]\,\Bigl(\underbrace{e^{-j\frac{2\pi}{N}}}_{W_N}\Bigr)^{kn} \qquad k=0,\ldots,N-1 \\
&= \sum_{n=0}^{N-1} x[n]\,W_N^{kn}
\end{align*}
\[ W_N = e^{-j\frac{2\pi}{N}} \]
\[ y = x * h \]
\[ y[n] = \sum_{\substack{k=0\\ N\ \text{tane}}} x[k]\,h[n-k] \]

$x[n]$ complex $\Rightarrow$ $N^{2}$ complex mult.

$N(N-1)$ summation

for $H$ \quad $N^{2}$ complex multiplication

$Y = X H \rightarrow N^{1}$

$\hookrightarrow y$ \qquad $N^{2} + N^{2} + N + N^{2} = 3N^{2} + N$

belli koşullar altında Fourier dönüşümü ilişkisi bulmak, daha verimli olabilir
\[ N = 2^{v} \]
\[
\begin{array}{r c c c c c}
 & x[0] & x[1] & x[2] & \cdots & x[N-1] \\
\text{even} & x[0] & x[2] & x[4] & \cdots & x[N-2] \\
\text{odd} & x[1] & x[3] & x[5] & \cdots & x[N-1]
\end{array}
\]
\begin{align*}
X[k] &= \sum_{n=0}^{N-1} x[n]\,W_N^{kn} \qquad k=0,\ldots,N-1 \\
&= \sum_{\substack{n=0\\ n\ \text{even}\\ \text{let } n=2r}}^{N-2} x[n]\,W_N^{kn} + \sum_{\substack{n=1\\ n\ \text{odd}\\ n=2r+1}}^{N-1} x[n]\,W_N^{kn}
\end{align*}
\[ n=2r+1 \Rightarrow r=\frac{n-1}{2} \qquad \begin{aligned} n=1 &\Rightarrow r=0 \\ n=N-1 &\Rightarrow r=\frac{N}{2}-1 \end{aligned} \]
\begin{align*}
&= \sum_{r=0}^{N/2-1} x[2r]\,\underbrace{W_N^{2rk}}_{W_{N/2}^{rk}} + \sum_{r=0}^{N/2-1} x[2r+1]\,\underbrace{W_N^{(2r+1)k}}_{\substack{W_N^{2rk}\,W_N^{k}\\ W_{N/2}^{rk}}}
\end{align*}
\[ r = 0,1,2,\ldots \]
\[ W_N = e^{-j\frac{2\pi}{N}} \]
\[ (W_N)^{2rk} = \left(e^{-j\frac{2\pi}{N}}\right)^{2rk} = \underset{W_{N/2}}{\kutula{e^{-j\frac{2\pi}{N/2}rk}}} = W_{N/2}^{rk} \]
\[ = \underbrace{\sum_{r=0}^{N/2-1} x[2r]\,W_{N/2}^{rk}}_{\frac{N}{2}\ \text{noktalı DFT}} + W_N^{k}\underbrace{\sum_{r=0}^{N/2-1} x[2r+1]\,W_{N/2}^{rk}}_{\frac{N}{2}\ \text{point DFT}} \]
\[ \underbrace{\frac{N}{2}\,\frac{N}{2} + \frac{N}{2}\,\frac{N}{2}}_{N^{2}/2} \]

\sayfa{6}
\vurgu[turuncu]{FFT}

$O(N\log_{2}N)$

\okunamadi{$N^{2}$}

\begin{center}
\begin{tabular}{rrr}
$N$ & $N^{2}$ & $N\log_{2}N$ \\ \hline
4 & 16 & 8 \\
16 & 256 & 64 \\
1024 & 1E6 & 10240
\end{tabular}
\end{center}
\end{document}
